\section*{Capítulo \ref{ch:b3:planetary-interiors-surfaces} --- Interiores y superficies planetarias}

\begin{solution}{exo:b3:planetary-interiors-surfaces:1}
$\tfrac52C\bar\rho = 2.5\times0.35\times5429 = 4750$; $x^2 = (4750 -
3300)/(5429 - 3300) = 0.681$, $x = 0.825$, un núcleo de $0.825\times2440 =
\qty{2010}{km}$; $\rho_c = 3300 + 2129/0.825^3 = \qty{7090}{kg.m^{-3}}$.
\end{solution}

\begin{solution}{exo:b3:planetary-interiors-surfaces:2}
$m = 3500\times\tfrac43\pi(25)^3 = \qty{2.29e8}{kg}$; $E = \tfrac12 m v^2 =
\tfrac12\times\num{2.29e8}\times(\num{2.0e4})^2 = \qty{4.6e16}{J}$, o
$\num{4.6e16}/\num{4.2e15} = \qty{11}{Mt}$. La regla da un cráter de
$20\times\qty{50}{m} = \qty{1}{km}$.
\end{solution}

\begin{solution}{exo:b3:planetary-interiors-surfaces:3}
$N_1 = 120/\num{100000} = \qty{1.2e-3}{km^{-2}}$, error relativo $1/\sqrt{120}
= \qty{9}{\percent}$. En la parte lineal, $T = N_1/c = \num{1.2e-3}/\num{4.875e-4}
= \qty{2.46}{Ga}$ (el término exponencial, $\num{8e-7}$, es despreciable), y
$\sigma_T = T/\sqrt{120} = \qty{0.22}{Ga}$.
\end{solution}

\begin{solution}{exo:b3:planetary-interiors-surfaces:4}
Área $4\pi(\num{1.8215e6})^2 = \qty{4.17e13}{m^2}$; potencia $2.5\times
\num{4.17e13} = \qty{1.04e14}{W} = \qty{104}{TW}$, 2.2 veces la de la Tierra
total, desde una superficie doce veces menor, y $2.5/0.092 = 27$ veces su
flujo.
\end{solution}

\begin{solution}{exo:b3:planetary-interiors-surfaces:5}
(a)~$\bar\rho = \num{6.417e23}/(\tfrac43\pi(\num{3.3895e6})^3) =
\qty{3934}{kg.m^{-3}}$; $g = GM/R^2 = \qty{3.73}{m.s^{-2}}$; $\rho gR =
\qty{4.97e10}{Pa}$. (b)~$19\mu/(2\rho gR) = 19\times\num{6.7e10}/\num{9.94e10}
= 12.8$, $k_2 = 1.5/13.8 = 0.109$. (c)~El valor medido es mayor en
2.6 desviaciones típicas: Marte da más que un planeta sólido con
la rigidez del manto, como ocurriría con un núcleo líquido en el interior.
\end{solution}

\begin{solution}{exo:b3:planetary-interiors-surfaces:6}
(a)~Luna: $4\times9.8/1.62 = \qty{24}{km}$ en roca cristalina,
\qty{12}{km} en sedimentos; Marte y Mercurio: $4\times9.8/3.7 =
\qty{10.6}{km}$ y \qty{5.3}{km}. Los \qty{20}{km} de la Luna quedan entre
sus dos predicciones: la regla se cumple con un error de unas pocas decenas por ciento. (b)~$Y =
\rho gD/3 = 2600\times9.8\times4000/3 = \qty{34}{MPa}$, la resistencia de
roca fracturada, no la de un cristal intacto.
\end{solution}

\begin{solution}{exo:b3:planetary-interiors-surfaces:7}
Marte: $\num{3.22e-14}(e^{6.93T} - 1) + \num{4.875e-4}\,T = \num{2.0e-3}$;
la estimación lineal $4.1$~Ga es demasiado antigua, siendo grande
allí; la bisección da $T = \qty{3.34}{Ga}$ (término exponencial
$\num{3.6e-4}$, término lineal $\num{1.63e-3}$). La función lunar da
\qty{2.39}{Ga}. Cada función está calibrada para un solo cuerpo: la tasa de
impactos y el tamaño de los cráteres que produce un proyectil dado difieren
entre la Luna y Marte.
\end{solution}

\begin{solution}{exo:b3:planetary-interiors-surfaces:8}
(a)~$3.551181/1.769138 = 2.0073$ y $7.154553/3.551181 = 2.0147$, cercanos a
$2$. (b)~$0.565247 - 0.844789 + 0.279542 = 0.000000$ por día, con una precisión de
$\num{1e-7}$: los tres movimientos medios están bloqueados, una resonancia que mantiene
las posiciones relativas de las lunas en cada conjunción. (c)~Los tirones se repiten en
los mismos puntos de la órbita, así que se suman en lugar de promediarse;
alimentan la excentricidad tan rápido como las mareas la amortiguan.
\end{solution}

\begin{solution}{exo:b3:planetary-interiors-surfaces:9}
$n = 2\pi/(1.769138\times86400) = \qty{4.111e-5}{s^{-1}}$, $nR =
\qty{74.9}{m.s^{-1}}$, $(nR)^5/G = \qty{3.53e19}{W}$;
$\tfrac{21}{2}e^2(nR)^5/G = 10.5\times\num{1.6e-5}\times\num{3.53e19} =
\qty{5.9e15}{W}$ por unidad de $k_2/Q$. Entonces $k_2/Q =
\num{1.04e14}/\num{5.9e15} = 0.018$.
\end{solution}

\begin{solution}{exo:b3:planetary-interiors-surfaces:10}
(a)~$x^2 = (\tfrac52C\bar\rho - \rho_m)/(\bar\rho - \rho_m)$ es lineal en
$C$, con pendiente $\tfrac52\bar\rho/(\bar\rho - \rho_m)$. (b)~La pendiente es
$2.5\times3340/52 = 161$: $x^2$ cambia en $0.161$ por cada $0.001$ de $C$.
Para $C = 0.395$, $x^2 = 0.036 + 0.161 = 0.197$, $x = 0.44$, $\rho_c =
\qty{3880}{kg.m^{-3}}$; para $C = 0.393$, $x^2 < 0$: no hay núcleo alguno.
(c)~Cerca de $C = \tfrac25$ el momento de inercia apenas nota un núcleo pequeño:
su tercera cifra significativa decide entre un núcleo grande y ligero y ninguno.
Se necesitan ondas sísmicas, mareas y el campo magnético.
\end{solution}

\begin{solution}{exo:b3:planetary-interiors-surfaces:11}
(a)~$\sigma_T = N_1/(10\,N_1')$: a \qty{1.0}{Ga}, $N_1 \approx N_1' \approx
\num{4.875e-4}$, $\sigma_T = \qty{0.10}{Ga}$; a \qty{3.8}{Ga}, $N_1 =
\num{1.07e-2}$ y $N_1' = \num{6.15e-2}$, $\sigma_T = \qty{0.017}{Ga}$,
seis veces menor. (b)~Con $N_1 = cT$, $N_1' = c$ y $\sigma_T =
cT/(\sqrt n\,c) = T/\sqrt n$. (c)~La propia función de cronología es
incierta: se apoya en unos pocos emplazamientos datados de la Luna y en el escalado
de la Luna a Marte, errores que ningún recuento reduce.
\end{solution}

\begin{solution}{exo:b3:planetary-interiors-surfaces:12}
(a)~La razón es $(n_ER_E/n_IR_I)^5(e_E/e_I)^2 = (0.4982\times0.8569)^5
\times(0.009/0.004)^2 = 0.0142\times5.06 = 0.072$: unos
\qty{7.5e12}{W}, o $\num{7.5e12}/(4\pi(\num{1.5608e6})^2) =
\qty{0.24}{W.m^{-2}}$. (b)~Un océano desacopla la corteza de hielo del
interior: la corteza se flexiona mucho más que una luna sólida, y su hielo
se disipa de forma distinta que la roca caliente de Ío, así que $k_2$ y $Q$ cambian ambos.
\end{solution}

\begin{solution}{pb:b3:planetary-interiors-surfaces:1}
\textbf{1.} $\bar\rho = \num{6.4169e23}/(\tfrac43\pi(\num{3.3895e6})^3) =
\qty{3934}{kg.m^{-3}}$.

\textbf{2.} $\tfrac52C\bar\rho = 3600$; $x^2 = (3600 - 3500)/(3934 - 3500) =
0.230$, $x = 0.479$, un núcleo de \qty{1620}{km}; $\rho_c = 3500 +
434/0.479^3 = \qty{7450}{kg.m^{-3}}$.

\textbf{3.} $x = 1830/3389.5 = 0.540$; $\rho_m = (3600 - 0.2915\times
3934)/(1 - 0.2915) = \qty{3460}{kg.m^{-3}}$, $\rho_c = 3462 +
472/0.1574 = \qty{6460}{kg.m^{-3}}$.

\textbf{4.} Es un tercio menor. La presión en el núcleo de Marte es mucho
inferior a la de la Tierra, así que su metal está menos comprimido; y un núcleo tan
ligero sugiere elementos mucho más ligeros mezclados con el hierro.

\textbf{5.} El \omterm{def:b3:planetary-interiors-surfaces:love-number}{número de Love} (2003) es muy anterior a InSight (2021): ya
demostró que al menos el núcleo externo es líquido y situó su
radio entre 1520 y \qty{1840}{km}, un intervalo que contiene el
valor sismológico.

\textbf{6.} $1 - (1 - 20/3389.5)^3 = \qty{1.8}{\percent}$ a $1 - (1 -
37/3389.5)^3 = \qty{3.2}{\percent}$ del volumen, y menos de la masa:
más ligero que el manto, reduce un poco la densidad efectiva del manto,
un efecto menor que la incertidumbre del modelo.

\textbf{7.} Desde \qty{1}{km}: 54; desde 1.41: 16; desde 2: 10; desde 2.83: 5;
desde 4, 5.66, 8 y 11.3: 2; desde 16: 1.

\textbf{8.} $N = 54/80000 = \qty{6.75e-4}{km^{-2}}$, $\sigma = \sqrt{54}/80000
= \qty{0.92e-4}{km^{-2}}$ (\qty{14}{\percent}).

\textbf{9.} Se esperan $54\times0.355 = 19$ cráteres de \qty{1.41}{km} y
$54\times0.147 = 7.9$ de \qty{2}{km}; se contaron 16 y 10, dentro de una
desviación típica ($\sqrt{19} = 4.4$, $\sqrt{7.9} = 2.8$): consistente.

\textbf{10.} No: $0.072/\num{6.75e-4} = 107$ veces por debajo. El recuento
registra el tiempo.

\textbf{11.} Los intervalos de igual anchura en $\log D$ cubren el rango de manera uniforme y
mantienen cantidades comparables de cráteres en los intervalos de una distribución empinada.

\textbf{12.} Son cráteres secundarios, formados por bloques lanzados desde un
cráter primario grande: no son impactos independientes procedentes del espacio, añadirían
cráteres que la función de cronología no cuenta.

\textbf{13.} $T = \num{6.75e-4}/\num{4.875e-4} = \qty{1.385}{Ga}$; el
término exponencial, $\num{3.22e-14}(e^{9.6} - 1) = \num{4.7e-10}$, es una
millonésima parte de $N_1$.

\textbf{14.} $\sigma_T = T/\sqrt{54} = \qty{0.19}{Ga}$: $1.38 \pm 0.19$~Ga.

\textbf{15.} $\num{6.75e-4}/\num{8.38e-4} = \qty{0.81}{Ga}$, más joven en
casi la mitad.

\textbf{16.} $31 + 12 + 7 + 2 + 0 + 2 + 2 = 56$ cráteres; $N(\qty{2}{km}) =
56/120000 = \qty{4.67e-4}{km^{-2}}$; $N_1 = \num{4.67e-4}/0.1473 =
\qty{3.17e-3}{km^{-2}}$. Bisección sobre $\num{3.22e-14}(e^{6.93T} - 1) +
\num{4.875e-4}\,T = \num{3.17e-3}$ da $T = \qty{3.54}{Ga}$ (el término
exponencial \num{1.45e-3}, término lineal \num{1.73e-3}).

\textbf{17.} $N_1' = \num{3.22e-14}\times6.93\,e^{6.93\times3.54} +
\num{4.875e-4} = \num{1.05e-2}$; $\sigma_T = \num{3.17e-3}/(\sqrt{56}\times
\num{1.05e-2}) = \qty{0.040}{Ga}$.

\textbf{18.} Llanura A: \qty{14}{\percent}; llanura B: \qty{1.1}{\percent}. En
la parte inicial, empinada, de la cronología, un cambio grande de la
densidad de cráteres significa un cambio pequeño de edad; en la parte lineal la edad es tan
incierta como el recuento.

\textbf{19.} $120000\times\num{3.17e-3}\times(1 - 0.1473) = 324$; el conjunto
lleva $48 + 8 = 56$.

\textbf{20.} $N_1 = 56/(120000\times0.8527) = \qty{5.47e-4}{km^{-2}}$,
$T_r = \num{5.47e-4}/\num{4.875e-4} = \qty{1.12}{Ga}$.

\textbf{21.} $\sigma = 1.12/\sqrt{56} = \qty{0.15}{Ga}$.

\textbf{22.} Un cráter de \qty{2}{km} tiene \qty{400}{m} de profundidad: la lava tenía
al menos ese espesor. Los cráteres mayores son más profundos y sus bordes se elevan más;
en parte quedaron rellenos y todavía se aprecian.

\textbf{23.} $1.38 \pm 0.19$ y $1.12 \pm 0.15$~Ga difieren en \qty{0.26}{Ga},
aproximadamente una desviación típica combinada ($\sqrt{0.19^2 + 0.15^2} = 0.24$):
los recuentos permiten un único episodio, aunque no lo demuestran.

\textbf{24.} La joven llanura de lava A tiene una edad modelo de
\textbf{1.4~mil~millones~de~años} ($1.38 \pm 0.19$~Ga).
\end{solution}
